\documentclass{article}
\usepackage[utf8]{inputenc}
\usepackage{geometry}
\geometry{
	a4paper,
	total={170mm,257mm},
	left=20mm,
	top=20mm,
}
\usepackage{graphicx}
\usepackage{titling}
\usepackage[american,RPvoltages]{circuiTikz}
\usepackage{tikz, tikz-cd}
\usepackage{pgfplots}
\usepackage{siunitx}
\usepackage{float}
\usepackage{enumitem}
\usepackage{amsmath,amsfonts,stmaryrd,amssymb} % Math packages
\usetikzlibrary{decorations.markings}
\title{Schmitt Trigger}
\author{Rodrigo Castillejos}
\date{July 2024}

\usepackage{fancyhdr}
\fancypagestyle{plain}{%  the preset of fancyhdr 
	\fancyhf{} % clear all header and footer fields
	\fancyfoot[L]{\thedate}
	\fancyhead[L]{Problem Resolution}
	\fancyhead[R]{\theauthor}
}
\makeatletter
\def\@maketitle{%
	\newpage
	\null
	\vskip 1em%
	\begin{center}%
		\let \footnote \thanks
		{\LARGE \@title \par}%
		\vskip 1em%
		%{\large \@date}%
	\end{center}%
	\par
	\vskip 1em}
\makeatother

\usepackage{lipsum}  
\usepackage{cmbright}

\begin{document}
	
	\maketitle
	
	\noindent\begin{tabular}{@{}ll}
		Autor & \theauthor\\

	\end{tabular}
	
	\section*{La realimentación positiva}
	La mayoría de los circuitos usan realimentación negativa, en este tipo de realimentación una señal de voltaje o corriente proporcional a la señal de salida es retornada a la terminal de entrada inversora del amplificador operacional. Sin embargo, la realimentgación positiva puede ser usada para realizar un número útil de funciones no lineales, por ejemplo el comparador schmitt trigger y los circuitos multivibradores
	
	\section*{El comparador Schmitt Trigger}
	A menudo, es útil comparar una señal de voltaje de entrada a un nivel de referencia conocido. Esto suele hacerse de manera electrónica usando un circuito comparador como el que se muetsra en la figura 1. La salida del comparador es un $1$ lógico cuando la señal de entrada excede el nivel de referencia y un $0$ lógico cuando la entrada es menor que el nivel de referencia. El comparador básico es simplemente un amplificador sin realimentación de muy alta ganancia. Cuando la señal de entrada excede el voltaje de referencia $V_{REF}$, la salida se satura a $V_{CC}$, cuando la señal de entrada es menor que $V_{REF}$ la salida se satura a $-V_{EE}$ como se indica en la característica de transferencia de la figura 1(b).

	\begin{figure}[H]
	\begin{tikzpicture}[scale=0.8]
		\ctikzset{bipoles/thickness=4}
		\ctikzset{monopoles/vcc/arrow={Stealth[width=6pt,length=9pt]}}
		\ctikzset{monopoles/vee/arrow={Stealth[width=6pt,length=9pt]}}
		\draw (0,0) node[op amp, noinv input up, anchor=+](OA){};
		\draw (OA.+) to[short] ++(-2,0)
		      to[V, invert, l=$v_I$] ++(0,-3) node[ground]{};
		\draw (OA.-) to[short] ++(-0.5,0)
		      to[short,  -o] ++(0,-1) node[below]{$V_{ref}$};
		\draw (OA.out) to[short, -o] ++(0.5,0) node[right]{$v_O$};
		\draw(OA.up) to[short] ++(0,0.4) node[vcc](VCC){$+V_{CC}$};
		\draw(OA.down) to[short] ++(0,-0.4) node[vee](VEE){$-V_{EE}$};
		\node[anchor=west] at (-1,-7){$(a) \, circuito$};
		\node[anchor=west] at (4,-7){$(b) \, caracteristica$};
		\node[anchor=west] at (4,-7.5){$de \, transferencia$};
		\node[anchor=west] at (11.5,-7){$(c) \, Respuesta \, del \, comparador$};
		\node[anchor=west] at (11.5,-7.5){$ante \, una \, senal \, de \, entrada \, con$};
		\node[anchor=west] at (11.5,-8){$ruido$};
			  
		\pgfplotsset{width=7cm,compat=1.18}
		\begin{axis}[xshift=5cm,
			 yshift=-2.8cm,
			 clip=false,
		     axis y line=left,
			 axis x line=middle,
			 xlabel=$v_I$,
			 ylabel=$v_O$, 
			 axis y line= center,
			 xmin=0,xmax=3.5,
			 ymin=-1.5,ymax=1.5,
			 xtick={1.5},
			 xticklabels={$\,\,\,\,\,\,\,\,\,\,\,\,\,\,\,\,\,V_{REF}$},
			 ytick={-1,1},
			 yticklabels={$-V_{EE}$,$V_{CC}$},
			 ]
			\addplot+[line width=1pt,mark=none,const plot]
			coordinates {(0,-1) (1.5,-1) (1.5,1) (2.5,1)};
		\end{axis}
		
		\pgfplotsset{width=7cm,compat=1.18}
		\begin{axis}[xshift=12cm,
			yshift=0cm,
			clip=false,
			axis y line=left,
			axis x line=middle,
			xlabel=$t$,
			ylabel=$v_I$, 
			axis y line= center,
			xmin=0,xmax=3.5,
			ymin=0,ymax=1.5,
			xtick={0},
			xticklabels={$$},
			ytick={0.8},
			yticklabels={$$},
			]
			\addplot+[line width=1pt,mark=none]
			coordinates {(0.1,0.1) (0.3,0.4) (0.4, 0.08) (0.5,0.6) (0.7, 0.2) (0.8, 0.6) (0.9, 0.3) (1,0.75) (1.1,0.45) (1.4, 1) (1.6, 0.7) (1.8, 1.1) (2.1, 0.7) (2.2, 1) (2.4, 0.9) (2.5, 1.1) (3, 1)};
			\addplot[color = black, dashed, thick] coordinates {(0,0.8) (3.4,0.8)};
			\addplot[color = black, dashed, thick] coordinates {(1.3,0.8) (1.3,-1)}; %vertical1
			\addplot[color = black, dashed, thick] coordinates {(1.53,0.8) (1.53,-1)}; %vertical2
			\addplot[color = black, dashed, thick] coordinates {(1.65,0.8) (1.65,-1)}; %vertical3
			\addplot[color = black, dashed, thick] coordinates {(2.02,0.8) (2.02,-1)}; %vertical4
			\addplot[color = black, dashed, thick] coordinates {(2.15,0.8) (2.15,-1)}; %vertical5
			\node[anchor=west] at (axis cs:2.2,0.6){$Noisy \, input$};
			\node[anchor=west] at (axis cs:2.2,0.5){$signal$};
		\end{axis}
		
		\pgfplotsset{width=7cm,compat=1.18}
		\begin{axis}[xshift=12cm,
			yshift=-6cm,
			clip=false,
			axis y line=left,
			axis x line=middle,
			xlabel=$v_I$,
			ylabel=$v_O$, 
			axis y line= center,
			xmin=0,xmax=3.5,
			ymin=-1.5,ymax=1.5,
			xtick={0},
			xticklabels={$$},
			ytick={-1,1},
			yticklabels={$V_{EE}$,$V_{CC}$},
			]
			\addplot+[line width=1pt,mark=none,const plot]
			coordinates {(0.5,-1) (1.3,-1) (1.3,1) (1.53,1) (1.53,-1) (1.65,-1) (1.65,1) (2.02, 1) (2.02,-1) (2.15,-1) (2.15,1) (3.3,1)};
			\addplot[color = black, dashed, thick] coordinates {(0,-1) (0.5,-1)}; %horizontal2
			\node[anchor=west] at (axis cs:-0.05,0.8){$comparator$};
			\node[anchor=west] at (axis cs:0.2,0.6){$output$};
		\end{axis}
	\end{tikzpicture}
	\caption{}	
	\label{fig:figura1}
	\end{figure}

Los amplificadores construidos para usarse como comparadores son diseñaods específicamente para ser capaces de saturarse a cualquiera de los dos voltajes extremos sin incurrir en retrasos internos de tiempo excesivos.. Sin embargo, ocurre un problema cuando los comparadores de alta velocidad son usados con señales que incluyen ruido como se indica en la figura 1(c), es decir, pueden ocurrir muchas transiciones conforme la señal de entrada cruza el nivel de referencia debido al ruido presente en la entrada. En sistemas digitales, por lo general se desea detecter ese umbral de manera limpia y generando sólo una punica transicióny el circuito \textbf{Schmitt-trigger} de la figura 2 ayuda a resolver este problema. 

\begin{figure}[H]
	\centering
	\begin{tikzpicture}[scale=0.8]
		\ctikzset{bipoles/thickness=4}
		\ctikzset{monopoles/vcc/arrow={Stealth[width=6pt,length=9pt]}}
		\ctikzset{monopoles/vee/arrow={Stealth[width=6pt,length=9pt]}}
		\draw (0,0) node[op amp, noinv input up, anchor=+](OA){};
		\draw (OA.+) to[short] ++(-2,0) 
		      to[short] ++(0,3) coordinate (p1) node[above]{$v_{REF}$};
		\draw (p1) to[R, l=$R_2$] (p1 -| OA.out) -- (OA.out);
		\draw (p1) to[R, *-, l_=$R_1$] ++(-4,0)
		      to[short] ++(0,-0.5) node[ground]{};
		\draw(OA.up) to[short] ++(0,0.4) node[vcc](VCC){$+V_{CC}$};
		\draw(OA.down) to[short] ++(0,-0.4) node[vee](VEE){$-V_{EE}$};
		\draw (OA.out) to[short, *-o] ++(1,0) node[right]{$v_O$};
		\draw (OA.-) to[short] ++(-2,0)
		      to[V, invert, l_=$v_I$] ++(0,-3) node[ground]{};
		      
		      
        \pgfplotsset{width=9cm,compat=1.18}
    \begin{axis}[xshift=7cm,
    	yshift=-2.8cm,
    	clip=false,
    	axis y line=left,
    	axis x line=middle,
    	xlabel=$v_I$,
    	ylabel=$v_O$, 
    	axis y line= center,
    	xmin=-2,xmax=2,
    	ymin=-1.5,ymax=1.5,
    	xtick={0},
    	xticklabels={$$},
    	ytick={-1},
    	yticklabels={$-V_{EE}$},
    	]
    	\addplot+[line width=1pt,const plot, mark=none]
    	coordinates {(-1,1) (0.5,-1) (1.5,-1)};
    	\draw[<-](0.55, 0.05)--(0.7, 0.3); %vo
    	\node[anchor=west] at (axis cs:0.6,0.4){$\beta V_{CC}$};
    	\addplot[color = black, dashed, thick] coordinates {(0,-1) (0.5,-1)}; %horizontal1
    	\addplot[color = black, dashed, thick] coordinates {(0.5,1) (1,1)}; %horizontal1
    	\node[anchor=west] at (axis cs:1,1){$V_{CC}$};
    \end{axis}
	\end{tikzpicture}
	\caption{}	
	\label{fig:figura2}
\end{figure}

El Schmitt Trigger usa un comparador cuyo voltaje de referencia se deriva a partir de la red divisora de voltaje a través de la salida. La señal de entrada se aplica en la terminal de entrada inversora y el voltaje de referencia en la terminal de entrada no inversora (realimentación positiva). Para voltajes de salida positivos, $V_{REF} = \beta V_{CC}$ y para voltajes de salida negativo $V_{REF} = -\beta V_{EE}$ donde 

\begin{align}
	\beta = \frac{R_1}{R_1 + R_2}
\end{align}
\end{document}